\section{Personal Thesis Syllabus}\label{personal-thesis-syllabus}

\subsection{Machine Learning for Regime-Conditional Return Prediction
(NEC /
MoE)}\label{machine-learning-for-regime-conditional-return-prediction-nec-moe}

\textbf{Prepared for:} Tom Koreslesjak\\
\textbf{Purpose:} personal learning plan before proposal and thesis
execution\\
\textbf{Proposal target:} September 2026\\
\textbf{First full draft target:} January 2027\\
\textbf{Created:} June 2026

\begin{center}\rule{0.5\linewidth}{0.5pt}\end{center}

\subsection{1. Thesis Direction}\label{thesis-direction}

The thesis should be framed as \textbf{machine learning work with
finance as the empirical domain}, not as a pure finance thesis. The
finance motivation is regime-dependent return predictability; the ML
contribution is the design, correction, training, and interpretation of
mixture-of-experts models for non-stationary financial prediction.

The current NEC code in \texttt{OP\ model/nec/nec\_hybrid\_explained.py}
is a useful starting point because it is clearly a rough
mixture-of-experts idea:

\begin{itemize}
\tightlist
\item
  \texttt{C\_L}: an LSTM sequence encoder.
\item
  \texttt{C\_F}: a 3-class softmax classifier head.
\item
  \texttt{NEC\_MLP}: a feedforward regression expert.
\item
  \texttt{NEC}: a wrapper that hard-routes predictions to either an
  extreme expert or a normal expert.
\end{itemize}

The key defect is central to the thesis: the classifier produces three
classes, but the routing rule only checks whether
\texttt{argmax(prob)\ ==\ 1}. Classes 0 and 2 are collapsed into the
same branch. So the model computes a 3-way gating distribution and then
uses only a binary decision. It also uses hard \texttt{argmax} routing,
so regression-loss gradients do not naturally train the classifier gate
end-to-end.

The corrected thesis version should turn this into a principled MoE:

\begin{Shaded}
\begin{Highlighting}[]
\NormalTok{p(z = k | x) = softmax(g\_k(x))}
\NormalTok{y\_hat = sum\_k p(z = k | x) f\_k(x)}
\end{Highlighting}
\end{Shaded}

where \texttt{g\_k} is the gate and \texttt{f\_k} are expert predictors.
The thesis can then compare soft routing, hard routing, sparse/top-k
routing, and classical regime-switching baselines.

\begin{center}\rule{0.5\linewidth}{0.5pt}\end{center}

\subsection{2. Critical Framing: What Is True, What Needs
Care}\label{critical-framing-what-is-true-what-needs-care}

\subsubsection{Strong part of the
argument}\label{strong-part-of-the-argument}

Regime-switching is a serious and recognizable finance tradition.
Hamilton-style Markov-switching models are canonical; asset-pricing and
portfolio-choice papers often ask whether return dynamics, volatility,
risk premia, or factor payoffs change across latent market states. This
gives the thesis a finance motivation that a committee can understand.

The clean thesis pitch is:

\begin{quote}
Classical regime-switching models represent state-dependent return
dynamics with latent Markov states and relatively restrictive
observation models. Modern MoE models replace the restrictive expert
models with flexible neural experts and replace or augment Markov gating
with feature-conditioned gating. This thesis tests whether modern MoE
methods improve regime-conditional return prediction and whether learned
expert assignments correspond to recognizable market regimes.
\end{quote}

\subsubsection{Point to avoid
overstating}\label{point-to-avoid-overstating}

MoE is not automatically a full generalization of Hamilton. Hamilton has
a latent state with \textbf{Markov persistence}. A standard MoE gate is
usually \textbf{feature-conditioned} and has no explicit transition
matrix. The defensible claim is:

\begin{itemize}
\tightlist
\item
  MoE generalizes the \textbf{expert / observation model} side.
\item
  MoE relaxes or replaces the \textbf{gating mechanism} with
  feature-conditioned routing.
\item
  A hybrid model can add Markov structure back if needed.
\end{itemize}

That distinction should appear in the thesis and proposal.

\subsubsection{Interpretability
hypothesis}\label{interpretability-hypothesis}

It is plausible that learned experts correspond to known regimes such as
high-volatility, low-volatility, bull, bear, recession, expansion, or
risk-on/risk-off. But this must be tested, not assumed. Experts may also
split by size, liquidity, sector, missingness, filing frequency, or
artifacts in the data.

The thesis should include explicit gate diagnostics:

\begin{itemize}
\tightlist
\item
  Expert utilization over time.
\item
  Gate entropy and collapse checks.
\item
  Correlation with VIX or realized market volatility.
\item
  Alignment with NBER recession dates, if macro data is used.
\item
  Expert exposure to market beta, size, sector, liquidity, and
  volatility.
\item
  Stability of expert behavior across walk-forward folds.
\end{itemize}

\begin{center}\rule{0.5\linewidth}{0.5pt}\end{center}

\subsection{3. Data Recommendation}\label{data-recommendation}

\subsubsection{What NEC actually needs}\label{what-nec-actually-needs}

NEC does not need ``SEC data'' specifically. It needs a supervised
panel:

\begin{Shaded}
\begin{Highlighting}[]
\NormalTok{(date, ticker, feature history) {-}\textgreater{} future return target}
\end{Highlighting}
\end{Shaded}

The existing code expects a sequence tensor:

\begin{Shaded}
\begin{Highlighting}[]
\NormalTok{x shape = (batch\_size, sequence\_length, input\_dim)}
\end{Highlighting}
\end{Shaded}

The LSTM gate reads recent history. The regressors use last-step
features. So the natural dataset is a \textbf{cross-sectional stock
panel} with a rolling history of price/volume/factor features and a
future-return label.

\subsubsection{Best first dataset: price-first, free
data}\label{best-first-dataset-price-first-free-data}

For a quant-trading path, the most relevant first version should be
price-first:

\begin{itemize}
\tightlist
\item
  Universe: large, liquid U.S. equities.
\item
  Frequency: daily.
\item
  History: ideally 10+ years if available.
\item
  Features: returns, volatility, volume, momentum, reversal, drawdown,
  market-relative features.
\item
  Target: future 5-day, 10-day, or 20-day return, preferably
  cross-sectionally ranked.
\end{itemize}

This is better than starting with SEC fundamentals because it gives
enough observations for ML and aligns more directly with trading-model
evaluation. SEC fundamentals update quarterly or annually, so they are
useful but sparse.

\subsubsection{Recommended free data
stack}\label{recommended-free-data-stack}

Use a staged data plan:

\begin{enumerate}
\def\labelenumi{\arabic{enumi}.}
\tightlist
\item
  \textbf{Stage A: synthetic OP pipeline data}

  \begin{itemize}
  \tightlist
  \item
    Purpose: verify the model, training loop, gate behavior, and
    metrics.
  \item
    Use before any real-data claim.
  \end{itemize}
\item
  \textbf{Stage B: free daily price data}

  \begin{itemize}
  \tightlist
  \item
    Use Stooq or another free daily-price source for U.S. equities.
  \item
    Use \texttt{yfinance} only as a prototyping fallback because it is
    unofficial.
  \item
    Include SPY or broad-market ETF data for market-state features.
  \end{itemize}
\item
  \textbf{Stage C: free regime/context data}

  \begin{itemize}
  \tightlist
  \item
    VIX history for volatility regimes.
  \item
    Kenneth French factor data for market, size, value, profitability,
    investment, momentum context.
  \item
    FRED macro series and recession indicators if needed.
  \end{itemize}
\item
  \textbf{Stage D: SEC fundamentals extension}

  \begin{itemize}
  \tightlist
  \item
    Use SEC EDGAR / companyfacts for quarterly and annual fundamentals.
  \item
    Add features such as revenue growth, earnings growth, leverage,
    book-to-market proxy, profitability, accruals, shares outstanding
    changes.
  \item
    Align by filing acceptance date, not fiscal period end, to avoid
    lookahead.
  \end{itemize}
\end{enumerate}

\subsubsection{Why SEC should not be the first price
source}\label{why-sec-should-not-be-the-first-price-source}

SEC filings are not a clean stock-price database. SEC is valuable for
fundamentals and filing metadata. Stock prices should come from a
market-data source. If the thesis says ``SEC data'', it should mean
SEC-derived fundamentals, not the main source of price history.

\subsubsection{Initial feature set}\label{initial-feature-set}

Start with features that are simple, defensible, and common in quant
work:

\begin{itemize}
\tightlist
\item
  Asset returns: 1-day, 5-day, 20-day, 60-day log returns.
\item
  Reversal: short-horizon return reversal, e.g.~negative 1-day or 5-day
  return.
\item
  Momentum: 20-day, 60-day, 120-day cumulative return.
\item
  Realized volatility: rolling 5-day, 20-day, 60-day standard deviation.
\item
  Downside volatility: rolling standard deviation of negative returns.
\item
  Drawdown: distance from recent rolling high.
\item
  Volume features: dollar volume, volume change, rolling volume z-score.
\item
  Market-relative features: asset return minus SPY return, asset
  volatility divided by market volatility.
\item
  Cross-sectional ranks: rank-normalized versions of key features each
  date.
\end{itemize}

\subsubsection{Initial target}\label{initial-target}

Use one main target and keep it stable:

\begin{Shaded}
\begin{Highlighting}[]
\NormalTok{y\_\{i,t\} = log(P\_\{i,t+h\} / P\_\{i,t\})}
\end{Highlighting}
\end{Shaded}

where \texttt{h} is 5, 10, or 20 trading days. For a thesis, 5-day and
20-day are the most reasonable starting horizons.

For evaluation, focus on:

\begin{itemize}
\tightlist
\item
  Cross-sectional rank-IC by date.
\item
  IC information ratio.
\item
  Decile long-short return.
\item
  Turnover and transaction-cost drag.
\item
  Sharpe only after costs and with skepticism.
\end{itemize}

Later, use residual returns or market-neutral returns:

\begin{Shaded}
\begin{Highlighting}[]
\NormalTok{residual\_return = asset\_return {-} beta * market\_return}
\end{Highlighting}
\end{Shaded}

This prevents the model from just learning market direction.

\begin{center}\rule{0.5\linewidth}{0.5pt}\end{center}

\subsection{4. Books and Core Resources}\label{books-and-core-resources}

\subsubsection{Core textbooks}\label{core-textbooks}

Each entry is tagged as either \textbf{Primary} (read the listed
chapters carefully) or \textbf{Reference} (consult only when the listed
topics come up).

\begin{enumerate}
\def\labelenumi{\arabic{enumi}.}
\tightlist
\item
  Hastie, Tibshirani, Friedman, \emph{The Elements of Statistical
  Learning}. --- \textbf{Primary}

  \begin{itemize}
  \tightlist
  \item
    Read: Ch 2 (already done), Ch 3, Ch 4, Ch 7, Ch 8 (§8.5 especially).
  \item
    Main use: statistical learning theory, linear models,
    regularization, model assessment, EM.
  \item
    Local file:
    \texttt{Books/The\ Elements\ of\ Statistical\ Learning.pdf}.
  \end{itemize}
\item
  Bishop, \emph{Pattern Recognition and Machine Learning}. ---
  \textbf{Primary}

  \begin{itemize}
  \tightlist
  \item
    Read: Ch 1 (§1.6 information theory), Ch 2 (§2.3 Gaussians), Ch 3
    (optional reinforcement), Ch 9 (mixture models and EM --- central),
    Ch 10 (§10.1--10.3 variational inference), Ch 13 (§13.1--13.3 HMMs).
  \item
    Main use: probability, Gaussian models, mixture models, EM,
    variational inference, graphical models, HMMs.
  \end{itemize}
\item
  Prince, \emph{Understanding Deep Learning}. --- \textbf{Primary}

  \begin{itemize}
  \tightlist
  \item
    Read: Ch 1--8 (deep-learning foundations), Ch 12 (RNN, LSTM, GRU).
  \item
    Skip: Ch 9 (CNNs), Ch 10 (transformers) unless added as later
    baselines.
  \item
    Main use: modern deep-learning foundations and PyTorch-compatible
    intuition.
  \end{itemize}
\item
  Goodfellow, Bengio, Courville, \emph{Deep Learning}. ---
  \textbf{Reference}

  \begin{itemize}
  \tightlist
  \item
    Consult only for: Ch 8 (deeper optimization theory: convergence,
    second-order methods, batch normalization); Ch 7 (additional
    regularization perspectives: adversarial training, data
    augmentation).
  \item
    Not read cover-to-relevant-chapters. Prince is the primary
    deep-learning reference.
  \item
    Local file: \texttt{Books/Deep+Learning+Ian+Goodfellow.pdf}.
  \end{itemize}
\item
  Lopez de Prado, \emph{Advances in Financial Machine Learning}. ---
  \textbf{Primary}

  \begin{itemize}
  \tightlist
  \item
    Read: Ch 2 (financial data structures), Ch 3 (labeling), Ch 4
    (sample weights), Ch 7 (cross-validation in finance).
  \item
    Main use: purged validation, embargoing, backtest overfitting,
    financial ML methodology.
  \item
    Local file:
    \texttt{Books/advances-in-financial-machine-learning-1nbsped-9781119482109-9781119482116-9781119482086.pdf}.
  \end{itemize}
\item
  Tsay, \emph{Analysis of Financial Time Series}. --- \textbf{Primary}

  \begin{itemize}
  \tightlist
  \item
    Read: Ch 4 (Markov-switching models --- especially §4.6). Skim Ch
    2--3 (AR/ARMA/GARCH) as background.
  \item
    Main use: financial time-series basics, volatility, Markov
    switching.
  \end{itemize}
\item
  Campbell, Lo, MacKinlay, \emph{The Econometrics of Financial Markets}.
  --- \textbf{Reference / selective}

  \begin{itemize}
  \tightlist
  \item
    Read: Ch 2 (§2.1--2.4 predictability of returns), Ch 5 (§5.1--5.3
    present-value relations).
  \item
    Not read cover-to-cover; selective for finance framing in the
    proposal and literature review.
  \end{itemize}
\item
  Bali, Engle, Murray, \emph{Empirical Asset Pricing}. ---
  \textbf{Primary}

  \begin{itemize}
  \tightlist
  \item
    Read: Ch 1 (§1.1--1.5 introduction), Ch 3 (§3.1--3.4 portfolio
    sorts), Ch 8 (§8.1--8.5 cross-sectional regressions and
    Fama-MacBeth).
  \item
    Main use: cross-sectional return prediction, portfolio sorting
    methodology, factor context.
  \end{itemize}
\end{enumerate}

\subsubsection{Papers to know before
proposal}\label{papers-to-know-before-proposal}

\begin{itemize}
\tightlist
\item
  Hamilton (1989), Markov-switching model.
\item
  Jacobs, Jordan, Nowlan, Hinton (1991), adaptive mixtures of local
  experts.
\item
  Shazeer et al.~(2017), sparsely-gated MoE and load balancing.
\item
  Fedus, Zoph, Shazeer (2022), Switch Transformer.
\item
  Ang and Bekaert (2002), regime switches in interest rates.
\item
  Daniel and Moskowitz (2016), momentum crashes.
\item
  Bailey and Lopez de Prado (2014), deflated Sharpe ratio.
\end{itemize}

\begin{center}\rule{0.5\linewidth}{0.5pt}\end{center}

\subsection{5. Learning Plan Overview}\label{learning-plan-overview}

The plan is split into two tracks:

\begin{enumerate}
\def\labelenumi{\arabic{enumi}.}
\tightlist
\item
  \textbf{Proposal-readiness track.} Enough theory, code, and finance
  context to write a serious proposal by September 2026.
\item
  \textbf{Thesis-execution track.} Deeper material and empirical work
  through January 2027.
\end{enumerate}

This is intentionally not locked to a fixed weekly hour count. If you
have more time, move faster. If you have less time, preserve the
milestone order.

Exercise numbers refer to the standard editions of the listed books. If
your edition differs, use the closest exercise in the same numbered
section.

\begin{center}\rule{0.5\linewidth}{0.5pt}\end{center}

\subsection{6. Part I: Programming
Foundations}\label{part-i-programming-foundations}

Part I is the programming literacy phase. The goal is fluency with
Python, NumPy, pandas, and basic plotting so that tooling never blocks
the rest of the syllabus. All mathematical content is introduced in
later parts, in dedicated \textbf{Math Foundations} sections placed
immediately before the modules that need that math --- this keeps math
grounded in upcoming applications rather than crammed all at the start.

\subsubsection{Module 1: Python, NumPy, pandas, and reproducible
notebooks}\label{module-1-python-numpy-pandas-and-reproducible-notebooks}

\textbf{Goal.} Be able to load data, build a panel, compute features,
train simple models, and make plots without fighting the tools.

\textbf{Topics.}

\begin{itemize}
\tightlist
\item
  NumPy arrays, broadcasting, vectorization.
\item
  pandas indexing, groupby, rolling windows, joins, missing data.
\item
  Matplotlib and seaborn for diagnostics.
\item
  Scikit-learn estimators, pipelines, train/test splits.
\item
  Jupyter notebooks versus scripts.
\item
  Reproducibility: random seeds, file paths, config objects.
\end{itemize}

\textbf{Readings by section number.}

\begin{itemize}
\tightlist
\item
  VanderPlas, \emph{Python Data Science Handbook}, Ch 2.1-2.7 for NumPy.
\item
  VanderPlas, Ch 3.1-3.11 for pandas.
\item
  VanderPlas, Ch 4.1-4.15 for plotting.
\item
  Scikit-learn user guide, sections 1.1, 3.1, 3.3, and 6.1.
\end{itemize}

\textbf{Focused exercise set.}

\begin{itemize}
\tightlist
\item
  VanderPlas Ch 2 notebook checks: reproduce the array slicing,
  broadcasting, and aggregation examples.
\item
  VanderPlas Ch 3 notebook checks: reproduce the groupby, merge, and
  time-series indexing examples.
\item
  Applied: build a \texttt{(date,\ ticker,\ features,\ target)} panel
  from daily prices.
\item
  Applied: recreate the OP pipeline's rank-IC calculation on synthetic
  data.
\end{itemize}

\textbf{Checkpoint.}

You can reproduce the data -\textgreater{} features -\textgreater{}
labels -\textgreater{} model -\textgreater{} metrics pipeline on
synthetic data.

\begin{center}\rule{0.5\linewidth}{0.5pt}\end{center}

\subsection{7. Part II: Linear Models and Classical
ML}\label{part-ii-linear-models-and-classical-ml}

Part II begins with two \textbf{Math Foundations} sections (A and B)
that teach the linear algebra and probability needed for everything
downstream. The two ML modules that follow (Module 4 on linear models
and Module 5 on model assessment) then use this math throughout.
Subsequent parts (III--V) will each introduce only the \emph{new} math
needed for their modules in smaller, focused Math Foundations sections.

\subsubsection{Math Foundations A: Linear algebra
essentials}\label{math-foundations-a-linear-algebra-essentials}

\textbf{Goal.} Understand the linear algebra behind regression, PCA,
embeddings, neural-network layers, and matrix calculus.

\textbf{Recommended ordering within this module.}

The module is best worked in four ordered passes, not in parallel. SVD
is the pivot --- almost everything downstream becomes clearer once SVD
is in your toolkit, and in particular ESL's ridge regression material
reads as just ``regularized least squares'' without SVD but as
``principled shrinkage of noise-dominated singular directions'' with it.

\begin{enumerate}
\def\labelenumi{\arabic{enumi}.}
\tightlist
\item
  \textbf{Strang foundations.} Work through Strang §2.3, 3.2--3.4,
  5.1--5.3, 6.1--6.3 (vector spaces, projections, orthogonality, least
  squares, eigenvalues and eigendecomposition). Do the Strang exercises
  from the focused exercise set as you go. This builds the prerequisites
  for SVD; you cannot do SVD cleanly without eigendecomposition first.
  The matrix-calculus identities you need for Modules 4 and 5 (gradient
  of a quadratic form, gradient of a norm) are derived in the focused
  exercise set below; the more general matrix-calculus reference
  (Petersen \& Pedersen's \emph{Matrix Cookbook}) is introduced later in
  Math Foundations C, where the deeper Jacobian and tensor-chain-rule
  machinery is actually load-bearing for neural networks.
\item
  \textbf{PCA and SVD focus block} (the dedicated subsection below).
  Take this as one unit in this position. The 3Blue1Brown videos first,
  then Strang's SVD chapter, then ESL §14.5, then Bishop §12.1--12.2. Do
  the focus-block exercises here.
\item
  \textbf{ESL Chapter 3 (§3.2--3.4).} With SVD now in your toolkit, work
  through linear regression, ridge, and lasso. The ridge derivation in
  the focused exercise set is the natural capstone here --- at this
  point you should be able to derive ridge in closed form and
  immediately interpret the result as SVD-basis shrinkage.
\item
  \textbf{Remaining focused exercises and the Module 2 checkpoint.} ESL
  Ch 3 problems, the MAP / Gaussian-prior view of ridge, and the overall
  module checkpoint to confirm everything is in place.
\end{enumerate}

\textbf{Readings by section number.}

\begin{itemize}
\tightlist
\item
  Strang, \emph{Linear Algebra and Its Applications}, Sec 2.3, 3.2-3.4,
  5.1-5.3, 6.1-6.3 (Step 1), plus the SVD chapter (Step 2, typically
  §6.7 or Ch 7 depending on edition).
\item
  ESL, Ch 3.2-3.4 as the applied regression use case for linear algebra
  (Step 3).
\item
  ESL §14.5 and Bishop §12.1--12.2 (Step 2, inside the PCA/SVD focus
  block).
\end{itemize}

The Matrix Cookbook is NOT a reading here. The Strang chapters above
give you the matrix-calculus identities Modules 4 and 5 actually require
(gradient of a quadratic form, projection geometry of least squares).
The full Matrix Cookbook is introduced in Math Foundations C, where the
Jacobian and tensor-chain-rule content becomes load-bearing for neural
networks.

\textbf{Topics.}

\begin{itemize}
\tightlist
\item
  Vector spaces, bases, rank, null space.
\item
  Projections and least squares.
\item
  Orthogonality and QR.
\item
  Eigenvalues and eigendecomposition.
\item
  SVD and PCA.
\item
  Positive semidefinite matrices.
\item
  Norms and conditioning.
\item
  Matrix calculus: gradients of quadratic forms, traces, matrix-vector
  products.
\item
  Jacobian of softmax.
\end{itemize}

\textbf{Focused exercise set.}

\emph{General linear algebra (Step 1 readings).}

\begin{itemize}
\tightlist
\item
  Strang Sec 3.3: Exercises 1, 3, and 7.
\item
  Strang Sec 6.3: Exercises 1 and 3.
\item
  Analytical: derive Gram--Schmidt orthogonalization for a 3-column
  matrix. Use it to derive the QR decomposition \(A = QR\) and show that
  the columns of \(Q\) form an orthonormal basis for the column space of
  \(A\).
\item
  Analytical: prove that a symmetric matrix \(A\) is positive
  semidefinite if and only if all its eigenvalues are non-negative. Then
  show that for any matrix \(X\), the Gram matrix \(X^\top X\) is always
  PSD. (You will use this fact constantly: it's why ridge regression's
  \(X^\top X + \lambda I\) is always invertible for \(\lambda > 0\).)
\item
  Analytical: derive the matrix-calculus identities
  \(\nabla_x (a^\top x) = a\),
  \(\nabla_x (x^\top A x) = (A + A^\top) x\) (and specialize to
  symmetric \(A\) to get \(2 A x\)), and
  \(\nabla_A \text{tr}(A^\top B) = B\). These are the workhorse
  identities you will use to derive almost every ML closed-form
  solution.
\item
  Analytical: derive the Jacobian of the softmax function
  \(\hat{p}_k = e^{z_k} / \sum_j e^{z_j}\) with respect to \(z\). Show
  \(\partial \hat{p}_k / \partial z_j = \hat{p}_k (\delta_{kj} - \hat{p}_j)\),
  so the Jacobian is \(\text{diag}(\hat{p}) - \hat{p}\hat{p}^\top\).
  (This Jacobian reappears as the Hessian of cross-entropy in Module 4.)
\item
  Analytical: define the condition number of a matrix as
  \(\kappa(A) = \sigma_{\max}/\sigma_{\min}\). Show that for ridge
  regression, the condition number of \(X^\top X + \lambda I\) is
  \((\sigma_{\max}^2 + \lambda)/(\sigma_{\min}^2 + \lambda)\), which
  approaches 1 as \(\lambda\) grows. This is the precise
  numerical-stability argument for regularization.
\end{itemize}

\emph{Linear regression and ridge (Step 3 readings).}

\begin{itemize}
\tightlist
\item
  ESL Ch 3: Exercises 3.4 and 3.6.
\item
  Analytical: derive the ridge regression closed-form solution
  \(\hat{\beta} = (X^\top X + \lambda I)^{-1} X^\top y\) from the
  regularized objective. Then, using the SVD \(X = U \Sigma V^\top\),
  show that ridge shrinks each singular direction by the factor
  \(\sigma_j^2 / (\sigma_j^2 + \lambda)\). This makes the SVD basis the
  natural ``eigenbasis of regularization.''
\item
  Analytical: derive ridge regression as MAP estimation under a Gaussian
  prior on \(\beta\); identify exactly which prior variance corresponds
  to a given \(\lambda\).
\item
  Analytical: show that the OLS solution
  \(\hat{\beta} = (X^\top X)^{-1} X^\top y\) is the orthogonal
  projection of \(y\) onto the column space of \(X\). Use this to derive
  the hat matrix \(H = X(X^\top X)^{-1} X^\top\) and verify it is
  symmetric and idempotent (\(H^2 = H\)).
\end{itemize}

\textbf{PCA and SVD focus block (3--4 days of dedicated work).}

You have not previously encountered PCA or SVD, so treat this as its own
learning thread within Module 2. The underlying math is the
eigendecomposition you already know from quantum mechanics (energy
eigenstates of a Hermitian Hamiltonian, \(H = U \Lambda U^\dagger\)),
but the applied technique, the statistical interpretation, and the
data-side intuition all need their own time.

\emph{Step 1: Intuition first (before any reading).} Watch 3Blue1Brown's
\emph{Essence of Linear Algebra} series on YouTube --- specifically the
eigenvalues/eigenvectors video and the SVD video (``But what is the
SVD?''). Twenty minutes total, exceptionally good visual intuition for
what these decompositions actually do to a point cloud and to a linear
map.

\emph{Step 2: Theory.}

\begin{itemize}
\tightlist
\item
  Strang, \emph{Linear Algebra and Its Applications}, the SVD chapter
  (typically §6.7 or Ch 7 depending on edition). Strang's treatment is
  direct and connects SVD to the four fundamental subspaces.
\item
  ESL §14.5 (Principal Components, Curves, and Surfaces). PCA in the
  unsupervised-learning context: when it works, when it fails, relation
  to factor analysis.
\item
  Bishop Ch 12 §12.1--12.2 (PCA and probabilistic PCA). PPCA recasts PCA
  as a maximum-likelihood estimator under a Gaussian latent-variable
  model --- the conceptual bridge to the mixture models you will see in
  Module 8.
\end{itemize}

\emph{Topics specific to this block.}

\begin{itemize}
\tightlist
\item
  SVD as the general decomposition \(A = U \Sigma V^\top\) for any
  matrix.
\item
  Relationship between SVD and eigendecomposition: they coincide for
  symmetric positive semidefinite matrices, differ otherwise.
\item
  PCA as eigendecomposition of the sample covariance matrix
  \(\hat{\Sigma} = X^\top X / n\).
\item
  PCA via SVD of the centered data matrix (numerically preferred).
\item
  Explained variance ratio and choosing the number of components.
\item
  Probabilistic PCA: PCA as maximum-likelihood estimation under a
  Gaussian latent-variable model with isotropic noise.
\item
  Connection to statistical factor models in finance (PCA-derived
  factors vs.~observed Fama-French factors; Connor and Korajczyk's
  APT-PCA tradition).
\item
  Why PCA can mislead in finance: signal-to-noise issues,
  non-stationarity of return covariances, and the random-matrix-theory
  cleaning argument (you'll come back to this if you use spectral
  methods later).
\end{itemize}

\emph{Focused exercises for the PCA/SVD block.}

\begin{itemize}
\tightlist
\item
  By hand, compute the SVD of a 2×2 or 3×3 matrix (e.g.,
  \(\begin{pmatrix} 3 & 1 \\ 1 & 3 \end{pmatrix}\) or another tractable
  example). Verify \(A = U \Sigma V^\top\) explicitly by multiplication.
\item
  Analytical: show that for a symmetric positive semidefinite matrix
  \(A\), the SVD reduces to eigendecomposition with \(U = V\) and
  \(\Sigma = \Lambda\) (the diagonal matrix of eigenvalues).
\item
  Analytical: derive PCA from a variance-maximization argument. Show
  that the first principal direction
  \(v_1 = \arg\max_{\|v\|=1} \text{Var}(X v)\) is the top eigenvector of
  the sample covariance matrix, and that subsequent directions are
  obtained by deflation.
\item
  Analytical: prove the equivalence of (a) PCA via eigendecomposition of
  \(X^\top X / n\) and (b) PCA via SVD of the centered data matrix
  \(X\). Show explicitly the relationship between the singular values
  and the eigenvalues.
\item
  Analytical: derive the maximum-likelihood solution of probabilistic
  PCA. Show that in the zero-noise limit \(\sigma^2 \to 0\), the PPCA
  solution recovers the standard PCA principal directions.
\item
  Connect to finance, on paper: explain why for a typical U.S. equity
  return matrix the first principal component is approximately the
  market factor --- i.e., a roughly equal-weighted average of stocks.
  Sketch what the second and third components typically capture (size
  and value-ish dimensions). This is an interpretation exercise, no
  computation needed.
\end{itemize}

\emph{Checkpoint for the PCA/SVD block.}

You can: explain SVD in one sentence; describe when SVD and
eigendecomposition coincide and when they differ; compute PCA two ways
(eigendecomposition of covariance vs.~SVD of centered data) and verify
they agree; interpret the explained variance ratio; explain PPCA as a
latent-variable model; explain in plain language why the first principal
component of a U.S. equity return matrix is ``the market.''

\textbf{Module 2 overall checkpoint.}

You can explain why ridge regression is a linear smoother and why
SVD/PCA matter for noisy financial features. You are fluent enough with
eigendecomposition, SVD, and PCA that they no longer require lookup ---
they are tools you reach for naturally when reasoning about covariance
structure, factor models, dimensionality reduction, and weight-matrix
analysis.

\subsubsection{Math Foundations B: Probability and statistics
essentials}\label{math-foundations-b-probability-and-statistics-essentials}

\textbf{Goal.} Learn the probability language needed for mixture models,
likelihoods, gates, and uncertainty.

\textbf{Readings by section number.}

\begin{itemize}
\tightlist
\item
  ESL Ch 2.1-2.9.
\item
  Bishop Ch 1.1-1.6.
\item
  Bishop Ch 2.1-2.3.
\end{itemize}

\textbf{Topics.}

\begin{itemize}
\tightlist
\item
  Random variables, expectation, variance, covariance.
\item
  Conditional probability and Bayes rule.
\item
  Conditional expectation.
\item
  Gaussian distributions, multivariate Gaussians.
\item
  Maximum likelihood and MAP.
\item
  Bias, variance, irreducible error.
\item
  Law of large numbers and central limit theorem intuition.
\item
  Bootstrap and sampling variability.
\item
  Heavy tails and why financial returns are not Gaussian.
\item
  KL divergence, entropy, cross entropy, mutual information.
\end{itemize}

\textbf{Focused exercise set.}

\emph{Textbook problems.}

\begin{itemize}
\tightlist
\item
  ESL Ch 2: Exercises 2.1, 2.3, 2.5, and 2.8.
\item
  Bishop Ch 1: Exercises 1.10 and 1.11.
\item
  Bishop Ch 2: Exercise 2.20.
\end{itemize}

\emph{MLE and Gaussians.}

\begin{itemize}
\tightlist
\item
  Analytical: derive the MLE for the mean and variance of a univariate
  Gaussian from \(n\) i.i.d. samples. Show that the MLE for variance is
  \(\hat{\sigma}^2_{\text{MLE}} = \tfrac{1}{n} \sum_i (x_i - \bar{x})^2\)
  and that this estimator is \emph{biased} --- specifically,
  \(\mathbb{E}[\hat{\sigma}^2_{\text{MLE}}] = \tfrac{n-1}{n} \sigma^2\).
  Derive the unbiased correction (Bessel's correction).
\item
  Analytical: derive the conditional distribution of a multivariate
  Gaussian. Given \(X = (X_1, X_2) \sim \mathcal{N}(\mu, \Sigma)\) with
  block structure, derive the conditional distribution
  \(p(X_1 \mid X_2 = x_2)\) in closed form --- both the conditional mean
  (the Schur complement formula) and the conditional covariance. (This
  formula appears constantly: it's how Kalman filters and Gaussian
  processes condition on observations.)
\item
  Analytical: derive the differential entropy
  \(h(\mathcal{N}(\mu, \sigma^2))\) of a univariate Gaussian. Then
  derive the maximum-entropy principle: show that among all
  distributions on \(\mathbb{R}\) with fixed mean \(\mu\) and variance
  \(\sigma^2\), the Gaussian uniquely maximizes the differential
  entropy. This is the deep reason Gaussians appear everywhere in
  statistics.
\end{itemize}

\emph{Bias-variance, high dimensions, KL.}

\begin{itemize}
\tightlist
\item
  Analytical: derive the bias-variance decomposition for the expected
  squared error of a generic estimator \(\hat{f}(x)\) of a target
  \(y = f(x) + \varepsilon\) with \(\mathbb{E}[\varepsilon] = 0\) and
  \(\text{Var}(\varepsilon) = \sigma^2\). Identify the three additive
  terms (irreducible noise, squared bias, variance) and verify the
  cross-terms vanish.
\item
  Analytical: prove the distance-concentration phenomenon in high
  dimensions. For two independent uniform random vectors in \([0,1]^d\),
  derive \(\mathbb{E}\|X - Y\|^2\) and \(\text{Var}(\|X - Y\|^2)\) as
  functions of \(d\), and show that the coefficient of variation
  \(\text{Var}(\|X - Y\|)^{1/2} / \mathbb{E}\|X - Y\|\) shrinks like
  \(1/\sqrt{d}\). Interpret what this means for k-NN in high dimensions.
\item
  Analytical: derive the KL divergence between two univariate Gaussians
  \(\mathcal{N}(\mu_1, \sigma_1^2)\) and
  \(\mathcal{N}(\mu_2, \sigma_2^2)\) in closed form.
\item
  Analytical: prove Jensen's inequality for a convex function
  \(\varphi\): \(\mathbb{E}[\varphi(X)] \geq \varphi(\mathbb{E}[X])\).
  Use it to prove that KL divergence is always non-negative:
  \(\text{KL}(p \| q) \geq 0\), with equality iff \(p = q\) almost
  everywhere. (This is the inequality underneath everything in Module
  8's ELBO derivations.)
\end{itemize}

\emph{Heavy tails (finance-relevant).}

\begin{itemize}
\tightlist
\item
  Analytical: derive the moments of a Student-\(t\) distribution with
  \(\nu\) degrees of freedom. Show that the \(k\)-th moment exists only
  for \(k < \nu\). Use this to explain why empirical financial returns,
  which look approximately Student-\(t\) with \(\nu \approx 3\)--\(5\),
  have well-defined means and variances but potentially undefined fourth
  moments (kurtosis).
\end{itemize}

\textbf{Checkpoint.}

You can explain why a MoE model is a conditional mixture model and why
its loss contains a log-sum-exp term.

\subsubsection{Module 4: Linear models, regularization, and
classification}\label{module-4-linear-models-regularization-and-classification}

\textbf{Readings by section number.}

\begin{itemize}
\tightlist
\item
  ESL Ch 3.1-3.4, 3.8.
\item
  ESL Ch 4.1-4.4.
\item
  Bishop Ch 3.1-3.3 as optional reinforcement.
\end{itemize}

\textbf{Topics.}

\begin{itemize}
\tightlist
\item
  Linear regression.
\item
  Ridge and lasso.
\item
  Logistic regression.
\item
  Multiclass softmax regression.
\item
  Regularization as constrained optimization.
\item
  Regularization as Bayesian prior.
\item
  Calibration and class probabilities.
\end{itemize}

\textbf{Focused exercise set.}

\emph{Textbook problems.}

\begin{itemize}
\tightlist
\item
  ESL Ch 3: Exercises 3.2, 3.4, and 3.6.
\item
  ESL Ch 4: Exercises 4.2 and 4.4.
\item
  Bishop Ch 3: Exercise 3.3.
\end{itemize}

\emph{Softmax and cross-entropy (the NEC gate math).}

\begin{itemize}
\tightlist
\item
  Analytical: derive the gradient of the multi-class cross-entropy loss
  \(\mathcal{L} = -\sum_k y_k \log \hat{p}_k\) with respect to the
  pre-softmax logits \(z_k\), where
  \(\hat{p}_k = e^{z_k} / \sum_j e^{z_j}\). Show that the gradient takes
  the clean form
  \(\partial \mathcal{L} / \partial z_k = \hat{p}_k - y_k\) (the
  prediction minus the one-hot target).
\item
  Analytical: derive the Hessian of the multi-class cross-entropy with
  respect to the logits,
  \(\partial^2 \mathcal{L} / \partial z_k \partial z_j\). Show it equals
  \(\text{diag}(\hat{p}) - \hat{p}\hat{p}^\top\) and prove this matrix
  is positive semidefinite, so the loss is convex in the logits.
\item
  Analytical: show that binary logistic regression is the \(K=2\)
  special case of multi-class softmax regression. Recover the standard
  logistic sigmoid form \(\sigma(z) = 1/(1 + e^{-z})\) from the softmax
  formula.
\end{itemize}

\emph{Regularization as constrained optimization and as Bayesian prior.}

\begin{itemize}
\tightlist
\item
  Analytical: write ridge and lasso as constrained optimization problems
  --- minimize the residual sum of squares subject to
  \(\|\beta\|_2^2 \leq t\) (ridge) or \(\|\beta\|_1 \leq t\) (lasso) ---
  and derive the KKT conditions. Show that the constrained form is
  equivalent to the Lagrangian form (unconstrained with penalty
  \(\lambda\)) with a one-to-one correspondence between \(t\) and
  \(\lambda\).
\item
  Analytical: derive the lasso solution coordinate-wise. For a single
  coordinate update, derive the soft-thresholding operator
  \(\hat{\beta}_j = \text{sign}(\tilde{\beta}_j) \cdot \max(|\tilde{\beta}_j| - \lambda, 0)\),
  where \(\tilde{\beta}_j\) is the OLS update. Contrast with ridge,
  which scales every coordinate by \(1/(1 + \lambda)\) but never sets
  coordinates to exactly zero. This is the math underneath ``lasso does
  feature selection, ridge does not.''
\item
  Analytical: derive lasso as MAP estimation under a Laplace
  (double-exponential) prior on \(\beta\). Show that the prior
  \(p(\beta) \propto e^{-|\beta|/b}\) leads to the L1 penalty in the
  negative log-posterior.
\end{itemize}

\textbf{Checkpoint.}

You can derive the softmax gradient and Hessian from scratch and explain
why the gradient form \(\hat{p} - y\) is what makes softmax regression
numerically well-behaved. You can derive both ridge and lasso as
Lagrangian solutions to constrained least-squares and as MAP estimators
under Gaussian and Laplace priors respectively. You understand the math
of the NEC gate at the level needed to debug it later.

\subsubsection{Module 5: Model assessment and financial
validation}\label{module-5-model-assessment-and-financial-validation}

\textbf{Readings by section number.}

\begin{itemize}
\tightlist
\item
  ESL Ch 7.1-7.12.
\item
  Lopez de Prado, AFML Ch 7.1-7.6.
\item
  Bailey and Lopez de Prado (2014), Sec 1-4.
\end{itemize}

\textbf{Topics.}

\begin{itemize}
\tightlist
\item
  Train/validation/test splits.
\item
  Cross-validation and why standard k-fold fails in finance.
\item
  Walk-forward validation.
\item
  Purging and embargoing.
\item
  Lookahead bias.
\item
  Survivorship bias.
\item
  Multiple testing and selection bias.
\item
  IC, rank-IC, ICIR.
\item
  Long-short decile backtests.
\item
  Transaction costs and turnover.
\end{itemize}

\textbf{Focused exercise set.}

\emph{Textbook problems.}

\begin{itemize}
\tightlist
\item
  ESL Ch 7: Exercises 7.1, 7.2, and 7.10.
\end{itemize}

\emph{Purging, embargoing, and leakage.}

\begin{itemize}
\tightlist
\item
  Analytical: describe the AFML purging and embargo algorithms
  (§7.4--7.5) in precise pseudocode, then identify and prove the leakage
  conditions they each eliminate. Construct a small counterexample where
  naive k-fold without purging leaks information from the test fold into
  the training fold via label overlap.
\item
  Analytical: characterize survivorship bias quantitatively. Set up a
  simple two-period model where survivors and non-survivors have
  systematically different return distributions; derive the bias in the
  apparent mean return when the analyst observes only the survivor
  sample. Show the bias is monotone in the survival hazard.
\end{itemize}

\emph{IC, rank-IC, and ICIR.}

\begin{itemize}
\tightlist
\item
  Analytical: derive the expectation and variance of date-level rank-IC
  (Spearman correlation between predictions and realized returns) under
  the null hypothesis of random predictions. Show that for a
  cross-section of \(N\) stocks, the standard error of the daily rank-IC
  scales like \(1/\sqrt{N-1}\).
\item
  Analytical: derive the IC information ratio (ICIR) as a
  \(t\)-statistic on the mean of the daily rank-IC time series.
  Specifically, ICIR \(= \bar{\text{IC}} / \text{SE}(\bar{\text{IC}})\)
  where \(\text{SE}(\bar{\text{IC}}) = \sigma_{\text{IC}}/\sqrt{T}\).
  Show how this relates to the out-of-sample Sharpe ratio of a strategy
  that scales positions by predicted IC.
\end{itemize}

\emph{Backtests, costs, and the Sharpe ratio.}

\begin{itemize}
\tightlist
\item
  Analytical: write the decile long-short portfolio return at date \(t\)
  explicitly as a function of (a) the predicted-rank assignment, (b) the
  realized returns, (c) the long and short leg weights, and (d) the
  turnover-adjusted transaction cost. Show how turnover enters as the L1
  norm of the weight change between consecutive dates.
\item
  Analytical: derive the deflated Sharpe ratio formula from Bailey and
  López de Prado (2014). Identify which terms correct for selection bias
  (the maximum over \(N\) trials problem) and which correct for
  non-normality of returns (skewness and kurtosis adjustments).
\end{itemize}

\emph{Multiple testing (the factor zoo problem).}

\begin{itemize}
\tightlist
\item
  Analytical: derive the Bonferroni correction for multiple hypothesis
  testing. Show that testing \(N\) hypotheses at level \(\alpha/N\) each
  controls the family-wise error rate (FWER) at \(\alpha\). Identify the
  conservatism of this correction when test statistics are correlated.
\item
  Analytical: derive the Benjamini--Hochberg procedure for false
  discovery rate (FDR) control. State precisely what FDR controls
  vs.~what FWER controls, and explain why FDR is often the more
  appropriate criterion for screening large numbers of trading signals.
\item
  Analytical: state Harvey, Liu \& Zhu's (2016) ``and the cross-section
  of expected returns'' multiple-testing critique: if hundreds of
  factors have been tested in the literature, what \(t\)-statistic
  threshold should we require for a newly proposed factor to be
  considered statistically significant? Derive a back-of-envelope answer
  using the Bonferroni-style argument.
\end{itemize}

\textbf{Checkpoint.}

You can defend why your evaluation protocol is not leaking future
information. You can articulate why multiple-testing corrections matter
for your thesis's claim that the corrected NEC beats baselines, and
explicitly state the family of tests over which your significance claim
is corrected.

\begin{center}\rule{0.5\linewidth}{0.5pt}\end{center}

\subsection{8. Part III: Deep Learning
Foundations}\label{part-iii-deep-learning-foundations}

Part III builds the neural-network machinery that Module 7 (sequence
models) and ultimately the corrected NEC will use. It opens with a
focused Math Foundations section on the optimization-theory and
matrix-calculus identities that backpropagation, SGD, Adam,
initialization theory, and batch normalization rely on. The linear
algebra and probability you already have from Part II are still
load-bearing here.

\subsubsection{Math Foundations C: Optimization theory and matrix
calculus for neural
networks}\label{math-foundations-c-optimization-theory-and-matrix-calculus-for-neural-networks}

\textbf{Goal.} Build the optimization and matrix-calculus machinery you
will need to derive backpropagation, analyze gradient descent's
convergence, understand momentum and Adam, and derive principled
initialization schemes. Most of the underlying calculus you already
have; the new content is its application to high-dimensional, non-convex
objectives.

\textbf{Readings by section number.}

\begin{itemize}
\tightlist
\item
  Boyd \& Vandenberghe, \emph{Convex Optimization}, Ch 1 (introduction),
  §3.1--3.2 (convex functions), §9.1--9.3 (gradient descent and Newton's
  method). Free PDF at stanford.edu/\textasciitilde boyd/cvxbook.
\item
  Goodfellow Ch 4 (numerical computation: overflow, conditioning,
  gradient-based optimization) and §8.1--8.3 (challenges in neural
  network optimization, ill-conditioning, saddle points).
\item
  Petersen \& Pedersen, \emph{Matrix Cookbook} --- this is the primary
  reference for the deep-learning matrix calculus you'll do here and in
  Module 8. Read §2.1 (derivatives of a determinant), §2.4 (derivatives
  of matrices and vectors), §2.5 (special matrix derivatives), §2.6
  (derivatives of traces), §2.7 (derivatives of vector norms). The
  Cookbook's §8 (Gaussian distributions, with explicit moment formulas)
  is a forward reference for Math Foundations D and Module 8's
  variational mean-field derivations --- skim now, return later. Free
  PDF at the Imm Technical University of Denmark site.
\end{itemize}

\textbf{Topics.}

\begin{itemize}
\tightlist
\item
  Convex vs.~non-convex functions; why neural network losses are
  non-convex.
\item
  Gradient descent: linear convergence for strongly convex objectives,
  \(O(1/k)\) for general convex.
\item
  Stochastic gradient descent and the variance of stochastic gradients.
\item
  Momentum and its interpretation as a damped second-order ODE.
\item
  Adaptive methods: Adam, RMSProp, AdaGrad, and bias correction.
\item
  Conditioning of the Hessian and its effect on gradient-descent speed.
\item
  Saddle points and the (more practical) issue of poorly-conditioned
  plateaus.
\item
  Matrix calculus for backpropagation: chain rule on tensor expressions,
  Jacobian-vector products.
\item
  Initialization theory: forward and backward variance preservation
  arguments.
\end{itemize}

\textbf{Focused exercise set.}

\begin{itemize}
\tightlist
\item
  Analytical: prove that for a strongly convex quadratic
  \(f(\beta) = \tfrac{1}{2} \beta^\top A \beta - b^\top \beta\), vanilla
  gradient descent with constant step size
  \(\eta < 2/\lambda_{\max}(A)\) achieves linear convergence:
  \(f(\beta_t) - f^* \leq (1 - \eta \lambda_{\min}(A))^t (f(\beta_0) - f^*)\).
\item
  Analytical: derive the condition number's role in gradient descent.
  Show that the number of iterations to reach
  \(\varepsilon\)-suboptimality scales as \(\kappa \log(1/\varepsilon)\)
  where \(\kappa = \lambda_{\max}/\lambda_{\min}\).
\item
  Analytical: derive the SGD variance term. Show that under unbiased
  stochastic gradients \(\mathbb{E}[\tilde{g}] = \nabla f\) with
  \(\text{Var}(\tilde{g}) = \sigma^2\), SGD converges in expectation to
  a neighborhood of the optimum whose size is proportional to
  \(\eta \sigma^2 / \lambda_{\min}(A)\). Identify the bias-variance
  tradeoff in step-size selection.
\item
  Analytical: derive momentum as the discretization of
  \(\ddot{\beta} + \gamma \dot{\beta} + \nabla f(\beta) = 0\) (a damped
  second-order ODE). Identify the friction \(\gamma\) and time step
  \(h\) corresponding to a given momentum coefficient
  \(\beta_{\text{momentum}}\) and learning rate \(\eta\).
\item
  Analytical: derive matrix chain rule for backprop. For
  \(\mathcal{L}(W) = g(Wx)\) where \(g: \mathbb{R}^m \to \mathbb{R}\)
  and \(W \in \mathbb{R}^{m \times n}\), show that
  \(\nabla_W \mathcal{L} = (\nabla_y g(y)) x^\top\) where \(y = Wx\).
  Generalize to deeper compositions.
\end{itemize}

\textbf{Checkpoint.}

You can analyze gradient-descent convergence, derive Adam from scratch,
and explain why initialization matters in terms of variance propagation.
You are equipped to derive backpropagation and analyze its failure modes
in Module 6.

\subsubsection{Module 6: Neural networks and
optimization}\label{module-6-neural-networks-and-optimization}

\textbf{Primary reading.}

\begin{itemize}
\tightlist
\item
  Prince Ch 1.1-1.4, 2.1-2.5, 3.1-3.6, 4.1-4.5, 5.1-5.6, 6.1-6.5,
  7.1-7.5, 8.1-8.5.
\end{itemize}

\textbf{Reference (consult as needed for specific topics).}

\begin{itemize}
\tightlist
\item
  Goodfellow Ch 8 for deeper optimization theory: convergence analysis,
  second-order methods, batch normalization mechanics.
\item
  Goodfellow Ch 7 for additional perspectives on regularization,
  especially adversarial training and data augmentation.
\end{itemize}

\textbf{Topics.}

\begin{itemize}
\tightlist
\item
  MLPs as function approximators.
\item
  Loss functions.
\item
  Backpropagation.
\item
  SGD, momentum, Adam.
\item
  Learning rates and schedulers.
\item
  Initialization.
\item
  Batch normalization.
\item
  Dropout.
\item
  Weight decay.
\item
  Early stopping.
\item
  Gradient clipping.
\end{itemize}

\textbf{Focused exercise set.}

\emph{Textbook problems.}

\begin{itemize}
\tightlist
\item
  Prince Ch 3: Problems 3.1 and 3.2.
\item
  Prince Ch 5: Problems 5.1 and 5.2.
\item
  Prince Ch 7: Problems 7.1 and 7.2.
\item
  Goodfellow Ch 6: Exercises 6.1 and 6.2.
\end{itemize}

\emph{Backpropagation.}

\begin{itemize}
\tightlist
\item
  Analytical: derive backpropagation explicitly for a 2-layer MLP
  \(\hat{y} = W_2 \, \sigma(W_1 x + b_1) + b_2\) with squared-error
  loss. Express \(\partial \mathcal{L} / \partial W_1\) and
  \(\partial \mathcal{L} / \partial W_2\) as products of forward-pass
  quantities and their Jacobians. Then generalize to \(L\) layers by
  induction.
\item
  Analytical: derive the cross-entropy + softmax composition's gradient.
  Starting from the result that
  \(\nabla_z \mathcal{L}_{\text{CE}}(\text{softmax}(z), y) = \hat{p} - y\),
  work backward through one linear layer to derive the gradient with
  respect to the layer weights. Verify this matches the structure of the
  standard implementation.
\end{itemize}

\emph{Optimization.}

\begin{itemize}
\tightlist
\item
  Analytical: derive the convergence rate of vanilla SGD for a strongly
  convex quadratic
  \(\mathcal{L}(\beta) = \tfrac{1}{2} \beta^\top A \beta - b^\top \beta\).
  Show that with constant step size \(\eta < 2/\lambda_{\max}(A)\), the
  expected suboptimality decreases by a factor
  \((1 - \eta \lambda_{\min}(A))^t\) per step.
\item
  Analytical: derive the Adam update with bias correction. Starting from
  the first-moment \(m_t = \beta_1 m_{t-1} + (1-\beta_1) g_t\) and
  second-moment \(v_t = \beta_2 v_{t-1} + (1-\beta_2) g_t^2\), derive
  the bias-corrected estimates \(\hat{m}_t = m_t / (1 - \beta_1^t)\) and
  \(\hat{v}_t = v_t / (1 - \beta_2^t)\). Show why bias correction
  matters in the early iterations (\(t\) small) and becomes negligible
  later.
\item
  Analytical: derive momentum as a discretization of a continuous-time
  ODE with friction. Identify the friction coefficient that corresponds
  to a given momentum hyperparameter \(\beta\) and step size \(\eta\).
\end{itemize}

\emph{Initialization and normalization.}

\begin{itemize}
\tightlist
\item
  Analytical: derive Xavier (Glorot) initialization for a
  fully-connected layer \(h = Wx\). For input variance
  \(\text{Var}(x_j) = \sigma_x^2\), find the weight initialization
  variance \(\text{Var}(W_{ij})\) such that the output variance equals
  the input variance (forward-pass preservation). Repeat for
  backward-pass preservation. Take the average for the standard Xavier
  formula.
\item
  Analytical: derive He initialization for a layer with ReLU activation.
  Account for the fact that ReLU zeros out half the inputs on average;
  show that this halves the effective variance and requires doubling the
  initialization variance: \(\text{Var}(W_{ij}) = 2/n_{\text{in}}\).
\item
  Analytical: derive the gradient of batch normalization. For a
  mini-batch \(x_1, \ldots, x_B\) normalized to
  \(\hat{x}_i = (x_i - \mu)/\sigma\) with \(\mu, \sigma\) computed from
  the batch, derive \(\partial \mathcal{L} / \partial x_i\) in terms of
  \(\partial \mathcal{L} / \partial \hat{x}_i\) and the batch
  statistics.
\end{itemize}

\emph{Practical milestone.}

\begin{itemize}
\tightlist
\item
  Applied: implement a 2-layer PyTorch MLP and verify gradients against
  \texttt{torch.autograd}. Compare your hand-derived gradients (from the
  backprop exercise above) numerically against \texttt{autograd} on a
  small input.
\end{itemize}

\textbf{Checkpoint.}

You can derive backprop, SGD convergence, Adam bias correction,
Xavier/He initialization, and the batch-norm gradient from scratch. You
can build, train, debug, and evaluate a small PyTorch model without
copying a training loop blindly.

\subsubsection{Module 7: Sequence models for NEC
gating}\label{module-7-sequence-models-for-nec-gating}

\textbf{Readings by section number.}

\begin{itemize}
\tightlist
\item
  Prince Ch 12.1-12.6.
\item
  Goodfellow Ch 10.1-10.10.
\item
  Olah, ``Understanding LSTM Networks,'' sections 1-5.
\end{itemize}

\textbf{Topics.}

\begin{itemize}
\tightlist
\item
  RNNs.
\item
  Backpropagation through time (BPTT) and truncated BPTT.
\item
  Vanishing and exploding gradients (the eigenvalue argument on the
  recurrent Jacobian).
\item
  GRU.
\item
  LSTM gates.
\item
  Sequence encoders versus sequence-to-sequence models.
\item
  Hidden state interpretation.
\item
  Why financial sequence models overfit easily.
\end{itemize}

\textbf{Focused exercise set.}

\emph{Textbook problems.}

\begin{itemize}
\tightlist
\item
  Prince Ch 12: Problems 12.1 and 12.2.
\item
  Goodfellow Ch 10: Exercises 10.1 and 10.2.
\end{itemize}

\emph{BPTT and vanishing gradients.}

\begin{itemize}
\tightlist
\item
  Analytical: derive backpropagation through time for a 3-step vanilla
  RNN with hidden update \(h_t = \tanh(W h_{t-1} + U x_t)\). Express the
  gradient \(\partial \mathcal{L}_T / \partial W\) as a sum over
  \(t = 1, \ldots, T\) of products of recurrent Jacobians
  \(\prod_{s=t+1}^{T} \partial h_s / \partial h_{s-1}\). Then generalize
  to length \(T\).
\item
  Analytical: prove the vanishing/exploding gradient theorem for vanilla
  RNNs. Let \(J_t = \partial h_t / \partial h_{t-1}\) and assume its
  spectral radius is \(\rho\). Show that \(\|\prod_{s} J_s\|\) grows or
  shrinks geometrically as \(\rho^T\), leading to exploding
  (\(\rho > 1\)) or vanishing (\(\rho < 1\)) gradients over long
  sequences. State the implication for training RNNs on long financial
  sequences.
\item
  Analytical: derive truncated BPTT and show that it introduces bias in
  the gradient estimate. Quantify the bias as a function of the
  truncation length \(k\) vs.~the true sequence length \(T\).
\end{itemize}

\emph{Gating mechanism analysis.}

\begin{itemize}
\tightlist
\item
  Analytical: derive the GRU cell's update equation
  \(h_t = (1 - z_t) \odot h_{t-1} + z_t \odot \tilde{h}_t\) from the
  reset gate \(r_t\), update gate \(z_t\), and candidate state
  \(\tilde{h}_t\). Show that when the update gate \(z_t \to 0\), the
  gradient of \(h_t\) with respect to \(h_{t-1}\) approaches the
  identity, suppressing vanishing gradients.
\item
  Analytical: derive the LSTM cell-state update
  \(c_t = f_t \odot c_{t-1} + i_t \odot g_t\) and show that when the
  forget gate \(f_t \to 1\), the cell-state gradient pathway is
  approximately the identity. This is the formal statement of ``LSTM
  gates protect gradient flow.''
\item
  Analytical: contrast LSTM and GRU formally. Identify which parameters
  and which gates are dropped in moving from LSTM to GRU. Argue at the
  level of parameter count and computational cost why GRU is often
  comparable or better than LSTM on small sequences.
\end{itemize}

\emph{Practical milestone.}

\begin{itemize}
\tightlist
\item
  Applied: implement a GRU cell from scratch in PyTorch and validate
  against \texttt{torch.nn.GRU}. Verify your forward-pass output matches
  and that gradients agree under autograd.
\item
  Applied (reading, not coding): re-read
  \texttt{OP\ model/nec/nec\_hybrid\_explained.py} and map every tensor
  shape. Identify exactly where the LSTM encoder hidden state \(h_T\) is
  consumed by the classifier head and how it interacts with the routing
  wrapper.
\end{itemize}

\textbf{Checkpoint.}

You can derive BPTT, the vanishing-gradient eigenvalue argument, the
GRU/LSTM gating gradient pathways, and explain exactly what
\texttt{C\_L} and \texttt{C\_F} are doing in the current NEC code at the
math level.

\begin{center}\rule{0.5\linewidth}{0.5pt}\end{center}

\subsection{9. Part IV: Mixture Models, HMMs, and
MoE}\label{part-iv-mixture-models-hmms-and-moe}

Part IV is the conceptual core of the thesis: mixture models, EM, HMMs,
and modern mixture-of-experts theory. Two focused Math Foundations
sections appear in this part --- D, on information theory and
variational analysis (needed for Module 8's ELBO), and E, on special
distributions and reparameterization (needed for Module 10's
Gumbel-Softmax and routing tricks).

\subsubsection{Math Foundations D: Information theory and variational
analysis}\label{math-foundations-d-information-theory-and-variational-analysis}

\textbf{Goal.} Deepen the information-theoretic and variational tools
you'll need to derive the ELBO, understand EM as coordinate ascent on a
free-energy functional, and reason about latent-variable models in
general.

\textbf{Readings by section number.}

\begin{itemize}
\tightlist
\item
  MacKay, \emph{Information Theory, Inference, and Learning Algorithms},
  Ch 2 (probability and inference) and Ch 4 (the source coding theorem).
  Free PDF at inference.org.uk.
\item
  Bishop §1.6 (information theory --- re-read with EM in mind) and §10.1
  (variational inference setup).
\item
  Murphy, \emph{Probabilistic Machine Learning: An Introduction}, §8.4
  (the EM algorithm) and §10.1 (variational inference) for an
  alternative perspective. Free chapters at probml.github.io.
\item
  Petersen \& Pedersen, \emph{Matrix Cookbook}, §8 (Gaussian
  distributions and explicit moment formulas) --- primary reference for
  the Gaussian-moment computations that appear in variational mean-field
  updates with Gaussian components. Used heavily in Bishop §10.2's
  variational mixture of Gaussians.
\end{itemize}

\textbf{Topics.}

\begin{itemize}
\tightlist
\item
  Entropy, joint entropy, conditional entropy, mutual information.
\item
  The chain rule of mutual information.
\item
  KL divergence as a Bregman divergence; its asymmetry and what it
  means.
\item
  Forward vs.~reverse KL: mean-seeking vs.~mode-seeking variational
  approximations.
\item
  Jensen's inequality and convex/concave function arguments.
\item
  The ELBO as Helmholtz free energy:
  \(-\mathcal{F}(q, \theta) = \mathbb{E}_q[\log p(X,Z|\theta)] + H(q)\).
\item
  Exponential families and conjugate priors as the natural setting for
  closed-form variational updates.
\item
  Reparameterization vs.~score-function gradient estimators (preview for
  Module 10).
\end{itemize}

\textbf{Focused exercise set.}

\begin{itemize}
\tightlist
\item
  Analytical: derive the chain rule of mutual information:
  \(I(X; Y, Z) = I(X; Y) + I(X; Z \mid Y)\). Use it to motivate the
  data-processing inequality \(I(X; Z) \leq I(X; Y)\) for
  \(X \to Y \to Z\) Markov chains.
\item
  Analytical: prove that for any joint distribution \(p(X, Z)\) and any
  variational distribution \(q(Z)\),
  \(\log p(X) = \mathcal{L}(q) + \text{KL}(q(Z) \| p(Z \mid X))\) where
  \(\mathcal{L}(q) = \mathbb{E}_q[\log p(X, Z) - \log q(Z)]\). This is
  the ELBO decomposition that underlies all variational methods.
\item
  Analytical: derive the variational update for mean-field
  \(q(Z) = \prod_i q_i(Z_i)\). Show that the optimal \(q_i^*\) given the
  other factors fixed is
  \(q_i^*(Z_i) \propto \exp(\mathbb{E}_{-i}[\log p(X, Z)])\). (This is
  Bishop's equation 10.9.)
\item
  Analytical: contrast forward KL \(\text{KL}(p \| q)\) with reverse KL
  \(\text{KL}(q \| p)\). Show that minimizing forward KL leads to a
  ``mean-seeking'' approximation (covers all modes of \(p\)) while
  reverse KL leads to ``mode-seeking'' (locks onto one mode). Connect to
  why VI typically uses reverse KL.
\item
  Analytical: state and prove that the Gaussian distribution maximizes
  entropy among all distributions with given mean and covariance. (You
  already derived the univariate case in Math Foundations B; generalize
  to multivariate.)
\end{itemize}

\textbf{Checkpoint.}

You can derive the ELBO three ways (Jensen's inequality, KL
decomposition, free energy), explain when forward vs.~reverse KL is
appropriate, and motivate why mean-field variational inference is
sometimes adequate and sometimes not. You are equipped to follow Module
8's variational EM derivation cold.

\subsubsection{Module 8: Mixture models, EM, and variational
inference}\label{module-8-mixture-models-em-and-variational-inference}

\textbf{Readings by section number.}

\begin{itemize}
\tightlist
\item
  Bishop Ch 9.1-9.4.
\item
  Bishop Ch 10.1-10.3 (variational inference, the evidence lower bound,
  variational mixture of Gaussians).
\item
  ESL Ch 8.5.
\end{itemize}

\textbf{Topics.}

\begin{itemize}
\tightlist
\item
  Latent-variable models.
\item
  Gaussian mixture models.
\item
  Responsibilities.
\item
  EM algorithm.
\item
  Log-sum-exp.
\item
  Local optima and initialization.
\item
  Mixture collapse.
\item
  Evidence lower bound (ELBO) and EM as ELBO maximization.
\item
  Variational inference for mixture models.
\item
  KL divergence between distributions on latent variables.
\item
  Connection to Bayesian neural networks and variational autoencoders
  (forward-looking optional context).
\end{itemize}

\textbf{Focused exercise set.}

\begin{itemize}
\tightlist
\item
  Bishop Ch 9: Exercises 9.7, 9.12, and 9.15.
\item
  Bishop Ch 10: Exercise 10.1 (ELBO derivation) and one of 10.5, 10.10,
  or 10.12 (variational mixture of Gaussians).
\item
  ESL Ch 8: Exercise 8.1 or 8.2.
\item
  Analytical: derive the M-step closed-form updates for a Gaussian
  mixture model from scratch. Show that the optimal mean is the
  responsibility-weighted average
  \(\mu_k = \sum_n \gamma(z_{nk}) x_n / N_k\) where
  \(N_k = \sum_n \gamma(z_{nk})\), derive the analogous formula for the
  covariance, and show the mixing coefficients become
  \(\pi_k = N_k / N\).
\item
  Analytical: prove the monotonicity property of EM --- show that each
  iteration of E-step + M-step monotonically increases (or leaves
  unchanged) the observed-data log-likelihood. Use the ELBO
  decomposition
  \(\log p(X) = \mathcal{L}(q, \theta) + \text{KL}(q \| p(Z|X))\).
\item
  Analytical: derive EM explicitly as coordinate ascent on the ELBO.
  Write out the E-step as \(q^* = p(Z | X, \theta^{\text{old}})\) and
  verify this closes the KL gap; then write the M-step as maximizing the
  expected complete-data log-likelihood under \(q^*\).
\item
  Analytical: characterize the principal failure modes of EM --- (a)
  component collapse (a Gaussian centered on a single data point with
  variance shrinking to zero), (b) identifiability under label
  permutation, (c) local-optimum dependence on initialization. For each,
  identify the structural property of the GMM likelihood that produces
  the failure mode.
\item
  Analytical: derive the log-sum-exp trick for numerical stability. Show
  that \(\log \sum_k e^{a_k} = a^* + \log \sum_k e^{a_k - a^*}\) where
  \(a^* = \max_k a_k\). Explain why this prevents overflow/underflow
  when computing mixture likelihoods
  \(\log \sum_k \pi_k \mathcal{N}(x \mid \mu_k, \Sigma_k)\) for inputs
  far from any component mean.
\item
  Analytical: derive the responsibility update of the E-step explicitly.
  Starting from Bayes' rule, show
  \(\gamma(z_{nk}) = p(z_n = k \mid x_n, \theta) = \pi_k \mathcal{N}(x_n \mid \mu_k, \Sigma_k) / \sum_j \pi_j \mathcal{N}(x_n \mid \mu_j, \Sigma_j)\).
  Verify this matches the variational \(q^*\) from the ELBO derivation.
\end{itemize}

\textbf{Checkpoint.}

You can write the MoE likelihood, explain how it relates to GMMs, state
EM as a coordinate ascent on the ELBO, derive M-step closed forms, apply
the log-sum-exp trick to prevent numerical issues, and articulate why
the variational view matters for any future Bayesian or stochastic-gate
extension of the MoE.

\subsubsection{Module 9: HMMs and classical regime
switching}\label{module-9-hmms-and-classical-regime-switching}

\textbf{Readings by section number.}

\begin{itemize}
\tightlist
\item
  Bishop Ch 13.1-13.3.
\item
  Tsay Ch 4.1-4.6, with Ch 4.6 as the Markov-switching focus.
\item
  Hamilton (1989), Sec 1-4.
\end{itemize}

\textbf{Topics.}

\begin{itemize}
\tightlist
\item
  Hidden Markov models.
\item
  Forward algorithm.
\item
  Backward algorithm.
\item
  Forward-backward smoothing.
\item
  Viterbi decoding.
\item
  Baum-Welch.
\item
  Hamilton filter.
\item
  Markov-switching mean and volatility.
\item
  Markov-switching regression.
\end{itemize}

\textbf{Focused exercise set.}

\begin{itemize}
\tightlist
\item
  Bishop Ch 13: Exercises 13.1, 13.3, and 13.5.
\item
  Tsay Ch 4: do two end-of-chapter problems tied to Markov switching or
  volatility regimes.
\item
  Analytical: derive the forward recursion for an HMM in full. Starting
  from \(\alpha_t(j) = p(o_1, \ldots, o_t, s_t = j)\), show that
  \(\alpha_t(j) = \big(\sum_i \alpha_{t-1}(i) a_{ij}\big) b_j(o_t)\),
  where \(a_{ij}\) is the transition probability and \(b_j\) the
  emission distribution. Compute the complexity: \(O(NT^2)\) via the
  forward recursion versus \(O(T^N)\) for a naive joint-distribution
  sum.
\item
  Analytical: derive the backward recursion \(\beta_t(i)\) symmetrically
  and show how forward-backward smoothing recovers the posterior
  marginals
  \(\gamma_t(i) = p(s_t = i \mid o_{1:T}) = \alpha_t(i) \beta_t(i) / p(o_{1:T})\).
\item
  Analytical: derive the Hamilton filter (1989) as the special case of
  HMM filtering applied to a Markov-switching mean and volatility model
  with Gaussian observations. Identify what changes between the general
  HMM filter and Hamilton's filter in terms of the observation model.
\item
  Analytical: state precisely and prove the structural difference
  between Hamilton-style Markov-switching gating (latent state has
  Markov persistence governed by a transition matrix \(P\)) and
  feature-conditioned MoE gating (latent state given by a deterministic
  function of features \(x_t\) with no transition structure). Show that
  the two reduce to the same model in the degenerate case where the
  transition matrix encodes only the marginal state distribution.
\item
  Analytical: derive the Viterbi algorithm as the max-product analog of
  the forward recursion. Define
  \(\delta_t(j) = \max_{s_1, \ldots, s_{t-1}} p(o_{1:t}, s_{1:t-1}, s_t = j)\)
  and show it satisfies the recursion
  \(\delta_t(j) = \big(\max_i \delta_{t-1}(i) a_{ij}\big) b_j(o_t)\).
  Use backpointer reconstruction to recover the most-likely state
  sequence \(\arg\max_{s_{1:T}} p(s_{1:T} \mid o_{1:T})\).
\item
  Analytical: derive the Baum-Welch EM updates for HMM parameters. Using
  the forward-backward posteriors
  \(\gamma_t(i) = p(s_t = i \mid o_{1:T})\) and pairwise marginals
  \(\xi_t(i,j) = p(s_t = i, s_{t+1} = j \mid o_{1:T})\), derive the
  M-step closed-form updates for the transition matrix
  \(\hat{a}_{ij} = \sum_t \xi_t(i,j) / \sum_t \gamma_t(i)\) and the
  emission parameters.
\end{itemize}

\textbf{Checkpoint.}

You can derive the full HMM machinery (forward, backward, Viterbi,
Baum-Welch) from scratch and state precisely how a Markov-switching
model differs from a feature-conditioned MoE.

\subsubsection{Math Foundations E: Special distributions and
reparameterization}\label{math-foundations-e-special-distributions-and-reparameterization}

\textbf{Goal.} Cover the discrete-distribution sampling and
reparameterization machinery you'll need for Module 10's Gumbel-Softmax,
straight-through estimator, and load-balancing derivations.

\textbf{Readings by section number.}

\begin{itemize}
\tightlist
\item
  Maddison, Mnih, Teh (2017), \emph{The Concrete Distribution: A
  Continuous Relaxation of Discrete Random Variables}, ICLR --- Sec
  1--3. Companion paper to Jang et al.~(2017) Gumbel-Softmax.
\item
  Kingma \& Welling (2014), \emph{Auto-Encoding Variational Bayes}, Sec
  2 --- for the original reparameterization trick.
\item
  Wikipedia article on the Gumbel distribution --- for the CDF, PDF, and
  max-stability property.
\end{itemize}

\textbf{Topics.}

\begin{itemize}
\tightlist
\item
  Extreme value distributions: Gumbel as the limiting distribution of
  the maximum of i.i.d. samples from a wide class of distributions
  (Fisher-Tippett theorem).
\item
  The Gumbel-Max trick: how Gumbel-perturbed logits sample exactly from
  a categorical distribution.
\item
  The Gumbel-Softmax relaxation: trading bias for differentiability via
  a temperature parameter.
\item
  Reparameterization trick more generally: when an expectation
  \(\mathbb{E}_{z \sim q_\phi}[f(z)]\) can be rewritten with a fixed
  base distribution as
  \(\mathbb{E}_{\varepsilon}[f(g_\phi(\varepsilon))]\) to enable
  gradient flow.
\item
  Score-function (REINFORCE) gradient estimator as the high-variance
  alternative for non-reparameterizable distributions.
\item
  Straight-through estimators: forward-pass discreteness with
  backward-pass continuity.
\end{itemize}

\textbf{Focused exercise set.}

\begin{itemize}
\tightlist
\item
  Analytical: derive the CDF and PDF of the standard Gumbel distribution
  from its definition \(G = -\log(-\log U)\) where
  \(U \sim \text{Uniform}(0, 1)\). Verify \(F_G(g) = e^{-e^{-g}}\).
\item
  Analytical: prove the Gumbel-Max trick. Let \(g_1, \ldots, g_K\) be
  logits and \(G_1, \ldots, G_K \sim \text{Gumbel}(0, 1)\) i.i.d. Show
  \(P(\arg\max_k (g_k + G_k) = j) = e^{g_j} / \sum_k e^{g_k}\). The
  proof uses the max-stability of the Gumbel distribution.
\item
  Analytical: derive the Gumbel-Softmax relaxation. Replace \(\arg\max\)
  with \(\text{softmax}_\tau\) at temperature \(\tau\), write
  \(y_k = e^{(g_k + G_k)/\tau} / \sum_j e^{(g_j + G_j)/\tau}\), and show
  that as \(\tau \to 0\), \(y \to\) one-hot at
  \(\arg\max_k(g_k + G_k)\). Identify the bias-variance tradeoff: small
  \(\tau\) means low bias (samples are nearly one-hot) but high gradient
  variance.
\item
  Analytical: derive the reparameterization gradient. For
  \(z \sim \mathcal{N}(\mu_\phi, \sigma_\phi^2)\), rewrite as
  \(z = \mu_\phi + \sigma_\phi \varepsilon\) with
  \(\varepsilon \sim \mathcal{N}(0, 1)\) fixed. Show that
  \(\nabla_\phi \mathbb{E}_z[f(z)] = \mathbb{E}_\varepsilon[\nabla_\phi f(\mu_\phi + \sigma_\phi \varepsilon)]\),
  swapping the gradient and expectation operators.
\item
  Analytical: derive the REINFORCE / score-function gradient estimator.
  For categorical \(z \sim \text{Cat}(\pi_\phi)\), show
  \(\nabla_\phi \mathbb{E}_z[f(z)] = \mathbb{E}_z[f(z) \nabla_\phi \log \pi_\phi(z)]\).
  Compare its variance to the reparameterized estimator and explain why
  Gumbel-Softmax is preferred when applicable.
\end{itemize}

\textbf{Checkpoint.}

You can derive the Gumbel-Max trick from scratch, explain when
Gumbel-Softmax vs.~REINFORCE is the right gradient estimator, and
articulate the role of temperature in the bias-variance tradeoff. You're
equipped to derive the routing-and-load-balancing machinery in Module
10.

\subsubsection{Module 10: Mixture-of-experts
theory}\label{module-10-mixture-of-experts-theory}

\textbf{Readings by section number.}

\begin{itemize}
\tightlist
\item
  Jacobs, Jordan, Nowlan, Hinton (1991), Sec 1-4.
\item
  Shazeer et al.~(2017), Sec 2-4.
\item
  Fedus, Zoph, Shazeer (2022), Sec 2-3.
\item
  Jang, Gu, Poole (2017), Sec 1-3.
\item
  Jiang et al.~(2024), \emph{Mixtral of Experts} --- read for the
  current state of the art in open-source MoE, with focus on top-2
  routing, expert-parallel training, and the auxiliary-loss choices used
  in practice.
\item
  Dai et al.~(2024), \emph{DeepSeekMoE} --- read for fine-grained expert
  specialization, shared expert isolation, and the load-balancing
  refinements that distinguish it from Switch / Mixtral.
\end{itemize}

\textbf{Topics.}

\begin{itemize}
\tightlist
\item
  Conditional mixture decomposition.
\item
  Soft routing.
\item
  Hard routing.
\item
  Sparse top-k routing.
\item
  Gumbel-softmax and straight-through estimators.
\item
  Load-balancing loss.
\item
  Expert collapse.
\item
  Gate entropy.
\item
  Expert specialization.
\item
  Mixture density networks.
\end{itemize}

\textbf{Focused exercise set.}

\emph{Paper-derivation exercises.}

\begin{itemize}
\tightlist
\item
  Paper exercise: reproduce the MoE objective from Jacobs et al.~Sec 2
  in your own notation. Show how it arises as a conditional mixture
  decomposition
  \(p(y \mid x) = \sum_k p(z = k \mid x) \, p(y \mid x, z = k)\) and
  identify what each factor parameterizes.
\item
  Paper exercise: reproduce the Shazeer auxiliary load-balancing loss
  from Sec 4 of Shazeer et al.~(2017). Derive it as
  \(\mathcal{L}_{\text{aux}} = \alpha \cdot K \cdot \sum_k f_k \cdot P_k\)
  where \(f_k\) is the fraction of inputs routed to expert \(k\) and
  \(P_k\) is the average gate probability. Show that this is minimized
  when both \(f_k\) and \(P_k\) are uniform across experts.
\end{itemize}

\emph{Routing and reparameterization.}

\begin{itemize}
\tightlist
\item
  Analytical: derive the Gumbel-Max trick. Let \(g_1, \ldots, g_K\) be
  gate logits and \(G_1, \ldots, G_K \sim \text{Gumbel}(0, 1)\) i.i.d.
  Prove that \(\arg\max_k (g_k + G_k)\) is distributed as a categorical
  sample with probabilities \(\hat{p}_k = e^{g_k} / \sum_j e^{g_j}\).
  (This requires deriving the Gumbel distribution's max statistic.)
\item
  Analytical: derive the Gumbel-Softmax relaxation as a continuous,
  differentiable approximation to argmax. Define
  \(y_k = e^{(g_k + G_k)/\tau} / \sum_j e^{(g_j + G_j)/\tau}\) and show
  that as \(\tau \to 0\), \(y \to\) one-hot at
  \(\arg\max_k (g_k + G_k)\). Explain the role of \(\tau\) as a
  bias-variance tradeoff for gradient estimation.
\item
  Analytical: derive the straight-through estimator. Show how it allows
  backpropagation through hard top-\(k\) routing by passing the gradient
  through as if routing were the identity, even though the forward pass
  returns a discrete selection.
\end{itemize}

\emph{Load balancing and expert dynamics.}

\begin{itemize}
\tightlist
\item
  Analytical: derive the gradient of the Shazeer auxiliary loss
  \(\mathcal{L}_{\text{aux}}\) with respect to the gating logits.
  Identify which term (\(f_k\) vs \(P_k\)) provides the gradient signal
  and which acts as a stop-gradient scaling factor.
\item
  Analytical: analyze the expert-collapse failure mode formally. Set up
  a small model with \(K = 2\) experts and a softmax gate; show that
  without any load-balancing penalty, the gradient dynamics of standard
  MSE training naturally concentrate all routing mass on whichever
  expert has the lower loss at initialization, regardless of whether the
  data is actually homogeneous. State the rate of collapse as a function
  of learning rate.
\end{itemize}

\emph{Mixture density connection.}

\begin{itemize}
\tightlist
\item
  Analytical: derive a mixture density network (MDN) training objective.
  Set up a model that outputs \(K\) mixture means \(\mu_k(x)\),
  variances \(\sigma_k^2(x)\), and mixing weights \(\pi_k(x)\), all as
  neural-network outputs. Write the negative log-likelihood loss
  \(-\log \sum_k \pi_k(x) \mathcal{N}(y \mid \mu_k(x), \sigma_k^2(x))\)
  and contrast its log-sum-exp structure with the soft-routed MoE loss
  for point prediction. Explain when MDN is the right choice
  (heteroscedastic, multimodal predictive distributions) and when
  MoE-for-prediction is sufficient.
\end{itemize}

\emph{Practical milestone (the thesis prototype).}

\begin{itemize}
\tightlist
\item
  Applied: implement a soft-routed MoE on synthetic regression data with
  built-in regime structure (\(y = f_1(x)\) when \(z = 0\),
  \(y = f_2(x)\) when \(z = 1\)). Verify the gate learns to recover the
  latent regime.
\item
  Applied: add load balancing and plot expert utilization \(f_k\) and
  gate entropy \(H(p(z \mid x))\) over training. Verify load balancing
  prevents the collapse failure mode analyzed above.
\end{itemize}

\textbf{Checkpoint.}

You can derive the modern MoE machinery from scratch --- soft routing,
Gumbel-Softmax, straight-through estimator, load balancing, and the
expert-collapse dynamics. You can build the corrected NEC prototype:

\begin{Shaded}
\begin{Highlighting}[]
\NormalTok{sequence features {-}\textgreater{} gate {-}\textgreater{} K experts {-}\textgreater{} weighted prediction}
\end{Highlighting}
\end{Shaded}

\begin{center}\rule{0.5\linewidth}{0.5pt}\end{center}

\subsection{10. Part V: Finance Domain Knowledge Needed for the
Thesis}\label{part-v-finance-domain-knowledge-needed-for-the-thesis}

Part V supplies the finance-domain knowledge the thesis sits in:
empirical asset pricing, factor models, and the practical
data-engineering realities of financial datasets. It opens with a brief
Math Foundations section on the econometric tools (time-series
vs.~cross-sectional regression, errors-in-variables, the structure of
factor models) that you'll use throughout Module 11 and the thesis's
evaluation chapter.

\subsubsection{Math Foundations F: Financial econometrics
primer}\label{math-foundations-f-financial-econometrics-primer}

\textbf{Goal.} Cover the regression-and-inference tools specific to
finance --- particularly the asymmetry between time-series and
cross-sectional regressions, the errors-in-variables structure of
two-pass estimation, and the GMM framing that underlies most empirical
asset pricing. Most of the math is OLS specialized to two different
sampling structures; the key insight is that \emph{which dimension you
sample over} changes which estimators are consistent.

\textbf{Readings by section number.}

\begin{itemize}
\tightlist
\item
  Cochrane, \emph{Asset Pricing}, Ch 11 (GMM in explicit detail) and Ch
  12 (regression-based tests). Sections 11.1--11.3 and 12.1--12.3 are
  the load-bearing pages. Free PDF available via Cochrane's University
  of Chicago page.
\item
  Campbell, Lo, MacKinlay Ch 5 (present-value relations) §5.1--5.3 for
  the basic econometric setup.
\item
  Greene, \emph{Econometric Analysis}, Ch 8 (asymptotic distributions of
  OLS) §8.1--8.3 --- re-read with finance-panel applications in mind.
\end{itemize}

\textbf{Topics.}

\begin{itemize}
\tightlist
\item
  Time-series regression: OLS of \(r_{i,t}\) on factors \(f_t\) over
  \(t = 1, \ldots, T\) for a fixed asset \(i\). Yields per-asset betas.
\item
  Cross-sectional regression: OLS of \(r_{i,t}\) on \(\hat{\beta}_i\)
  over \(i = 1, \ldots, N\) for a fixed (or average) date. Yields factor
  risk premia.
\item
  Why the two-pass (Fama-MacBeth) procedure is needed: time-series-only
  or cross-section-only fails to estimate risk premia consistently.
\item
  Errors-in-variables: the second-pass regressor \(\hat{\beta}_i\) is
  itself an estimate, with sampling error that biases the risk-premium
  estimate toward zero.
\item
  The Shanken correction: closed-form adjustment for errors-in-variables
  in the second-pass standard errors.
\item
  GMM as the unifying framework: factor models as moment conditions
  \(\mathbb{E}[(r_{i,t} - \alpha_i - \beta_i f_t) f_t] = 0\).
\item
  Heteroskedasticity and serial correlation in financial residuals;
  Newey-West standard errors.
\item
  The cross-section vs.~time-series \(R^2\) distinction.
\end{itemize}

\textbf{Focused exercise set.}

\begin{itemize}
\tightlist
\item
  Analytical: derive the OLS estimator for a single factor regression
  \(r_{i,t} = \alpha_i + \beta_i f_t + \varepsilon_{i,t}\). Show
  \(\hat{\beta}_i = \text{Cov}(r_i, f) / \text{Var}(f)\). Derive its
  asymptotic distribution under standard regularity conditions.
\item
  Analytical: derive the Fama-MacBeth two-pass estimator formally. State
  the regularity conditions under which the second-pass risk-premium
  estimator \(\hat{\lambda}\) converges to the true \(\lambda\) as
  \(T \to \infty\) for fixed \(N\).
\item
  Analytical: derive the errors-in-variables bias in the second-pass
  regression. Show that when \(\hat{\beta}_i = \beta_i + u_i\) with
  \(u_i\) measurement noise, OLS of returns on \(\hat{\beta}_i\)
  produces a slope estimate biased toward zero. Identify the attenuation
  factor in terms of
  \(\text{Var}(\beta) / (\text{Var}(\beta) + \text{Var}(u))\).
\item
  Analytical: derive the GMM objective for a factor model. Write the
  moment conditions \(\mathbb{E}[g_t(\theta)] = 0\) where
  \(g_t(\theta) = (r_{i,t} - \alpha_i - \beta_i f_t) \otimes (1, f_t^\top)^\top\).
  Show how minimizing the sample-moment quadratic form recovers OLS with
  the appropriate weighting matrix.
\item
  Analytical: derive Newey-West standard errors as a kernel-smoothed
  estimate of the long-run variance of residuals. Show why ordinary OLS
  standard errors are too small in the presence of serial correlation in
  returns.
\end{itemize}

\textbf{Checkpoint.}

You can derive the Fama-MacBeth estimator and its errors-in-variables
correction, write a factor model as a GMM problem, and explain why
Newey-West standard errors are the standard choice for financial panels.
You're equipped to engage with Module 11's empirical asset-pricing
material from first principles rather than as a list of techniques.

\subsubsection{Module 11: Empirical asset pricing for ML
people}\label{module-11-empirical-asset-pricing-for-ml-people}

\textbf{Readings by section number.}

\begin{itemize}
\tightlist
\item
  Bali, Engle, Murray Ch 1.1-1.5, Ch 3.1-3.4, Ch 8.1-8.5.
\item
  Campbell, Lo, MacKinlay Ch 2.1-2.4 and Ch 5.1-5.3.
\item
  Fama-French factor documentation, sections ``Daily Factors'' and
  ``Momentum Factor''.
\item
  Daniel and Moskowitz (2016), Sec 1-5.
\end{itemize}

\textbf{Topics.}

\begin{itemize}
\tightlist
\item
  Cross-sectional versus time-series prediction.
\item
  Factor models.
\item
  Market beta.
\item
  Size, value, momentum, profitability, investment.
\item
  Return predictability and low signal-to-noise.
\item
  Factor crashes and conditional performance.
\item
  Fama-MacBeth regression.
\item
  Residual returns and neutralization.
\end{itemize}

\textbf{Focused exercise set.}

\emph{Textbook problems.}

\begin{itemize}
\tightlist
\item
  Bali, Engle, Murray Ch 3: Exercises 1 and 2.
\item
  Campbell, Lo, MacKinlay Ch 2: Exercises 1 and 2.
\end{itemize}

\emph{Factor models and Fama-MacBeth.}

\begin{itemize}
\tightlist
\item
  Analytical: derive the Fama-MacBeth two-pass procedure from first
  principles. In the first pass, run time-series regressions of asset
  returns on factors to estimate per-asset factor betas. In the second
  pass, run cross-sectional regressions of returns on the estimated
  betas to estimate factor risk premia. State the assumptions under
  which the two-pass estimator is consistent and identify the
  errors-in-variables problem introduced by using estimated betas as
  regressors.
\item
  Analytical: derive the formula for the cross-sectional R-squared of a
  factor model and explain why time-series R-squared and cross-sectional
  R-squared can diverge.
\item
  Analytical: derive the OLS estimator for market beta
  \(\hat{\beta}_i = \text{Cov}(r_i, r_m) / \text{Var}(r_m)\) from a
  time-series regression
  \(r_{i,t} = \alpha_i + \beta_i r_{m,t} + \varepsilon_{i,t}\). Show how
  using stale (window-lagged) beta estimates avoids look-ahead bias when
  the beta is itself an input to a forecast.
\end{itemize}

\emph{Residual returns and neutralization.}

\begin{itemize}
\tightlist
\item
  Analytical: derive the formula for the market-residual return
  \(r_i^{\text{residual}} = r_i - \hat{\beta}_i r_m\). Show that the
  residual has zero correlation with the market by construction. Explain
  why residualizing eliminates market-factor exposure but does not
  eliminate cross-asset correlations driven by other (non-market)
  factors.
\item
  Analytical: derive multi-factor residualization. Given a set of \(K\)
  factors \(f_1, \ldots, f_K\), derive the residualized return as the
  residual from a multivariate regression of \(r_i\) on the factors.
  Show that the residualized cross-section is orthogonal to the factor
  space.
\end{itemize}

\emph{Factor crashes.}

\begin{itemize}
\tightlist
\item
  Analytical: state Daniel and Moskowitz's (2016) momentum-crash
  mechanism informally and then formally. Show that when the
  past-winners portfolio loaded short the market in the prior period
  (bear-market lookback), the momentum strategy implicitly takes a short
  market position right when the market is rebounding. Identify the
  conditional-beta structure that produces the crash.
\end{itemize}

\emph{Practical milestones (data-required).}

\begin{itemize}
\tightlist
\item
  Applied: download Fama-French factors from Kenneth French's data
  library and compute the market excess return series. This is
  operationally required for the thesis baseline.
\item
  Applied: run a cross-sectional regression of future returns on
  momentum and volatility features on a small panel. This sets up the
  methodology you will use for the thesis evaluation.
\end{itemize}

\textbf{Checkpoint.}

You can explain why a quant fund usually cares about cross-sectional
ranking, not only time-series price prediction.

\subsubsection{Module 12: Financial data
engineering}\label{module-12-financial-data-engineering}

\textbf{Topics.}

\begin{itemize}
\tightlist
\item
  Ticker universes.
\item
  Delistings and survivorship bias.
\item
  Corporate actions and adjusted prices.
\item
  Splits and dividends.
\item
  Point-in-time fundamentals.
\item
  Filing-date alignment.
\item
  Missing data.
\item
  Liquidity filters.
\item
  Rebalancing frequency.
\end{itemize}

\textbf{Readings by section number.}

\begin{itemize}
\tightlist
\item
  Lopez de Prado, AFML Ch 2.1-2.5, Ch 3.1-3.4, Ch 4.1-4.4.
\item
  SEC EDGAR API documentation, sections ``Submissions'' and ``Company
  Facts''.
\item
  SEC fair-access guidance, sections 1-3.
\end{itemize}

\textbf{Focused exercise set.}

\begin{itemize}
\tightlist
\item
  Build a large-cap ticker universe.
\item
  Pull daily price data for the universe and verify adjusted prices.
\item
  Build rolling features without lookahead.
\item
  Pull SEC companyfacts for 10 tickers and align by filing date.
\end{itemize}

\textbf{Checkpoint.}

You have a real-data panel ready for a baseline model.

\begin{center}\rule{0.5\linewidth}{0.5pt}\end{center}

\subsection{11. Proposal-Readiness
Milestones}\label{proposal-readiness-milestones}

Before writing the September proposal, have these artifacts:

\begin{enumerate}
\def\labelenumi{\arabic{enumi}.}
\tightlist
\item
  \textbf{Corrected NEC on synthetic data}

  \begin{itemize}
  \tightlist
  \item
    Soft-routed MoE with K = 2 or K = 3.
  \item
    Gate utilization and entropy plots.
  \item
    Synthetic hidden-regime recovery test.
  \end{itemize}
\item
  \textbf{Classical baseline}

  \begin{itemize}
  \tightlist
  \item
    Single MLP or LightGBM model.
  \item
    Simple HMM or Hamilton-style regime baseline.
  \end{itemize}
\item
  \textbf{Real price-data baseline}

  \begin{itemize}
  \tightlist
  \item
    Free daily-price panel.
  \item
    Simple technical/factor features.
  \item
    Forward-return target.
  \item
    Rank-IC and decile backtest.
  \end{itemize}
\item
  \textbf{Evaluation protocol}

  \begin{itemize}
  \tightlist
  \item
    Walk-forward validation.
  \item
    Purging and embargoing.
  \item
    Transaction costs.
  \item
    No random shuffling across time.
  \end{itemize}
\item
  \textbf{Literature map}

  \begin{itemize}
  \tightlist
  \item
    Classical regime switching.
  \item
    Classical MoE.
  \item
    Modern MoE.
  \item
    Financial ML validation.
  \item
    Cross-sectional return prediction.
  \end{itemize}
\item
  \textbf{One-page thesis claim}

  \begin{itemize}
  \tightlist
  \item
    What exact ML method is being tested?
  \item
    What finance prediction problem is used?
  \item
    What are the baselines?
  \item
    What would count as success?
  \item
    What would a negative result still teach?
  \end{itemize}
\end{enumerate}

\begin{center}\rule{0.5\linewidth}{0.5pt}\end{center}

\subsection{12. Thesis Research Design}\label{thesis-research-design}

\subsubsection{Main research question}\label{main-research-question}

Do modern mixture-of-experts routing methods improve out-of-sample
cross-sectional return prediction relative to single-model baselines and
classical regime-switching baselines?

\subsubsection{Secondary research
question}\label{secondary-research-question}

Do learned expert assignments correspond to recognizable financial
regimes or do they split the data along other dimensions?

\subsubsection{Candidate models}\label{candidate-models}

Start simple and expand only after the basics work:

\begin{enumerate}
\def\labelenumi{\arabic{enumi}.}
\tightlist
\item
  Ridge or elastic-net regression.
\item
  LightGBM.
\item
  Single MLP.
\item
  Original NEC hard-routed model.
\item
  Corrected soft-routed NEC / MoE.
\item
  Sparse top-k MoE.
\item
  HMM / Markov-switching regression baseline.
\end{enumerate}

\subsubsection{Main evaluation metrics}\label{main-evaluation-metrics}

\begin{itemize}
\tightlist
\item
  Mean daily rank-IC.
\item
  Rank-IC information ratio.
\item
  Decile long-short return.
\item
  Turnover.
\item
  Transaction-cost-adjusted Sharpe.
\item
  Maximum drawdown.
\item
  Expert utilization.
\item
  Gate entropy.
\item
  Gate-regime alignment.
\end{itemize}

\subsubsection{Minimum ablations}\label{minimum-ablations}

\begin{itemize}
\tightlist
\item
  Number of experts: K = 2, 3, 4, 8.
\item
  Soft versus hard routing.
\item
  With versus without load balancing.
\item
  Price-only features versus price + market-context features.
\item
  Raw returns versus market-residual returns.
\item
  Different prediction horizons: 5-day and 20-day.
\end{itemize}

\subsubsection{What a good negative result looks
like}\label{what-a-good-negative-result-looks-like}

If MoE does not beat strong baselines, the thesis can still be strong if
it shows:

\begin{itemize}
\tightlist
\item
  Single-model baselines are hard to beat in noisy financial data.
\item
  Expert gates collapse without regularization.
\item
  Learned experts do not align with intuitive regimes.
\item
  Apparent improvements disappear after purged validation or transaction
  costs.
\item
  Classical regime labels explain little about cross-sectional returns
  at the chosen horizon.
\end{itemize}

That is still a useful ML-for-finance result.

\begin{center}\rule{0.5\linewidth}{0.5pt}\end{center}

\subsection{13. Thesis Execution
Timeline}\label{thesis-execution-timeline}

\subsubsection{June to July 2026: foundations and code
fluency}\label{june-to-july-2026-foundations-and-code-fluency}

\begin{itemize}
\tightlist
\item
  Finish Python/pandas/NumPy/PyTorch foundations.
\item
  Work through ESL Ch 2, Ch 3, Ch 4, Ch 7.
\item
  Implement linear/logistic regression and MLP from scratch.
\item
  Understand the existing NEC code completely.
\item
  Run the OP pipeline tests and demos.
\end{itemize}

\subsubsection{July to August 2026: mixtures, HMMs, and corrected
NEC}\label{july-to-august-2026-mixtures-hmms-and-corrected-nec}

\begin{itemize}
\tightlist
\item
  Finish Bishop mixture-model material.
\item
  Implement GMM and EM.
\item
  Implement small HMM/forward algorithm.
\item
  Implement soft-routed MoE on synthetic data.
\item
  Add load-balancing and gate diagnostics.
\item
  Write the first corrected NEC prototype.
\end{itemize}

\subsubsection{August to September 2026: real-data baseline and
proposal}\label{august-to-september-2026-real-data-baseline-and-proposal}

\begin{itemize}
\tightlist
\item
  Build free daily-price panel.
\item
  Build simple price/volume features.
\item
  Train ridge, LightGBM, MLP baselines.
\item
  Add HMM/regime baseline.
\item
  Run walk-forward evaluation.
\item
  Write proposal with realistic scope.
\end{itemize}

\subsubsection{September to October 2026: thesis model
experiments}\label{september-to-october-2026-thesis-model-experiments}

\begin{itemize}
\tightlist
\item
  Train corrected NEC/MoE variants.
\item
  Compare soft, hard, and sparse routing.
\item
  Tune number of experts.
\item
  Add gate diagnostics.
\item
  Decide whether SEC fundamentals are feasible as an extension.
\end{itemize}

\subsubsection{October to November 2026: robust
evaluation}\label{october-to-november-2026-robust-evaluation}

\begin{itemize}
\tightlist
\item
  Run full walk-forward experiments.
\item
  Add transaction costs and turnover.
\item
  Run ablations.
\item
  Compare raw and residual targets.
\item
  Stress-test across horizons and regimes.
\end{itemize}

\subsubsection{November to December 2026: interpretation and
writing}\label{november-to-december-2026-interpretation-and-writing}

\begin{itemize}
\tightlist
\item
  Interpret expert behavior.
\item
  Compare gate assignments to VIX, realized volatility, SPY drawdowns,
  and recession indicators if used.
\item
  Write methods, data, and results chapters.
\end{itemize}

\subsubsection{December 2026 to January 2027: full
draft}\label{december-2026-to-january-2027-full-draft}

\begin{itemize}
\tightlist
\item
  Finish literature review.
\item
  Finish results tables and figures.
\item
  Write limitations.
\item
  Write conclusion.
\item
  Prepare first full draft.
\end{itemize}

\begin{center}\rule{0.5\linewidth}{0.5pt}\end{center}

\subsection{14. Core Exercise Spine}\label{core-exercise-spine}

If time gets compressed, preserve these:

\begin{enumerate}
\def\labelenumi{\arabic{enumi}.}
\tightlist
\item
  ESL Ch 2: Exercises 2.1, 2.5, and one bias-variance simulation.
\item
  ESL Ch 3: Exercises 3.4 and 3.6, plus ridge regression from scratch.
\item
  ESL Ch 4: Exercise 4.2, plus softmax regression with a gradient check.
\item
  Bishop Ch 9: Exercises 9.7 and 9.12, plus a GMM implementation.
\item
  Bishop Ch 13: Exercises 13.1 and 13.3, plus the HMM forward algorithm.
\item
  Prince Ch 7: Problems 7.1 and 7.2, plus a PyTorch MLP gradient check.
\item
  Shazeer et al.~Sec 4: reproduce the load-balancing loss and implement
  it.
\item
  AFML Ch 7.4-7.5: implement purging and embargoing.
\item
  Build a real cross-sectional price panel and compute rank-IC.
\item
  Compare corrected NEC against a single MLP and LightGBM baseline.
\end{enumerate}

\begin{center}\rule{0.5\linewidth}{0.5pt}\end{center}

\subsection{15. Final Personal Rule}\label{final-personal-rule}

Do not let the thesis become ``I used a complicated neural network on
stocks.'' The defensible thesis is:

\begin{Shaded}
\begin{Highlighting}[]
\NormalTok{I studied conditional computation for non{-}stationary financial prediction.}
\NormalTok{I corrected a flawed hard{-}routed NEC architecture into a principled MoE.}
\NormalTok{I compared routing mechanisms under proper financial validation.}
\NormalTok{I tested whether learned experts correspond to market regimes.}
\end{Highlighting}
\end{Shaded}

That is ML-first, finance-relevant, and realistic for a January 2027
first draft.

\begin{center}\rule{0.5\linewidth}{0.5pt}\end{center}

\subsection{16. Useful Source Links}\label{useful-source-links}

\begin{itemize}
\tightlist
\item
  SEC EDGAR APIs and data:
  https://www.sec.gov/search-filings/edgar-application-programming-interfaces
\item
  SEC fair access guidance: https://www.sec.gov/os/accessing-edgar-data
\item
  Stooq historical data: https://stooq.com/
\item
  Kenneth French data library:
  https://mba.tuck.dartmouth.edu/pages/faculty/ken.french/data\_library.html
\item
  FRED API: https://fred.stlouisfed.org/docs/api/fred/
\item
  CBOE VIX historical data:
  https://www.cboe.com/tradable\_products/vix/vix\_historical\_data/
\item
  Shazeer et al.~MoE paper page: https://research.google/pubs/pub45929/
\item
  Switch Transformer JMLR: https://www.jmlr.org/papers/v23/21-0998.html
\end{itemize}
